\justifying
\NSFCBodyText

\subsubsection{年度研究计划}

\textbf{第一阶段（立项后第1--12个月）。} 围绕“新能源数据分析与功率预测”示范课程模块的4个知识模块，完成公开、授权或自主编写资源的来源登记、授权审查和脱敏检查；建立产业任务—课程知识轻量结构与资源单元描述模式，形成课程文本、设备图/结构图和运行曲线/图表三类核心模态的解析流程。实验代码和产业案例仅在许可与格式条件满足时纳入扩展测试。构建人工整理、通用大模型直接构建和文本检索增强等基线，制定资源单元边界、知识点覆盖、专业属性和来源证据的标注规范。完成小规模开发集上的联合知识约束原型与预实验，冻结主要指标、数据划分和专家盲评量表。

阶段性验收以材料合规清单、资源元数据卡、知识结构与标注规范、开发集划分和可复现基线为主。对授权尚未明确、来源不能追溯或存在敏感风险的材料实行退出机制，不以扩张样本量为由降低合规要求；对预实验中未表现出稳定作用的约束，保留其失败记录并在第二阶段通过消融进一步判断是否应简化。

\textbf{第二阶段（立项后第13--24个月）。} 完善专业属性与来源约束的跨模态关联方法，以及内容、课程、模态和证据四层质量门控；在按知识模块、设备类型和材料来源隔离的测试集上完成对照、消融和跨来源留出验证。组织独立专家盲评和误差分析，形成可追溯的资源单元、资源包、元数据与质量报告；根据自动门控和人工修订结果迭代本地原型。项目实施期按两年安排，具体时间节点以立项通知书为准。

阶段性验收以冻结的测试配置、B0--B3 比较、模块消融、分层误差分析、专家盲评一致性和资源质量报告为主。若跨来源验证或专家盲评未支持预设机制，将形成边界与负结果说明，并将原型降级为“自动初筛—人工复核”工具；不将未达到验证条件的资源单元纳入最终资源包。

\subsubsection{预期研究结果}

预期形成以下成果：

\begin{enumerate}
  \item 产业任务—课程知识联合约束的多模态资源单元构建方法，以及资源边界、知识覆盖和专业属性的评价口径；
  \item 专业属性与来源证据约束的跨模态关联方法，形成可追溯关系标注和错误类型分析；
  \item 内容、课程、模态、证据四层资源内在质量控制方法，并给出与独立专家盲评相对应的验证报告；
  \item 面向示范课程模块的可追溯多模态资源包、资源元数据规范和离线或本地原型；
  \item 力争发表与项目内容直接相关的二区 SCI 论文2篇；其他成果仅在满足实验室任务书和实际研究进展时申报，不作预先承诺。
\end{enumerate}

成果交付遵循可复核原则：资源包中每一资源单元保留课程知识、专业属性、关联关系、来源位置和质量状态；方法报告保留数据划分、基线与消融设置、指标定义和误差类型；对外展示仅使用公开、授权或自主编写材料，且不含学生个人信息、成绩、学习行为或企业敏感信息。

\subsubsection{学术交流与合作计划}

项目拟在资源质量量表、专业术语映射和成果验证等环节，邀请具有新能源专业教学或相关工程背景的教师、专家开展同行讨论和盲评；在不涉及未授权材料、学生数据和企业敏感数据的前提下，参加人工智能、多模态学习、能源工程教育等相关学术会议或学术活动，交流资源构建和质量验证方法。合作以方法讨论、资源授权核验和专家评价为主，不以接入教育云或高校生产系统为前提。
